\FloatBarrier
\section{FT inner shield bolts}

The sixth beamline study renames nut1--nut9 to FTInnerShieldAndFlange\_Bolt1 through
FTInnerShieldAndFlange\_Bolt9. These engineering meshes are solid copper bolts, despite their former names.
The approved simplification removes threads and replaces the hexagonal head details with round heads and
smooth solid shafts. Each surface uses 24 angular sides, an annular shoulder, and closed end caps: 192
triangles per bolt. No hollow bore is introduced. G4\_Cu, colors, mothers, and passive placements are
retained across ten applicable CAD variations and ninety GEMC2 SQLite rows; only their names change in the
placement records.

The shaft axes, head shoulders, and outer ends are measured from each original mesh. The head radius is
14.2875 mm, within the original hexagonal across-flats envelope. The original tip caps contact the shield
blind-hole end faces. Shortening each inner tip by 0.15 mm creates an explicit gap; the smooth shaft radius
is then solved from the original signed solid volume. The seven regular shafts have radius approximately
8.5741 mm; the two longer shafts, bolts 4 and 8, have radius approximately 8.3681 mm. Their individual
float32 STL volumes agree with the originals within 0.0001\%. This redistributes copper locally and changes
external surfaces rather than preserving the engineering threaded profile.

Combined triangles fall from 16,278 to 1,728 (89.38\% fewer), and bytes fall from 814,656 to 87,156.
Original and optimized total copper volumes are 150623.172389 and 150623.196214 mm$^3$, a change of
0.00001582\%. All nine outputs are closed, singly connected, manifold genus-zero surfaces with no detected
self-intersections or zero-area triangles. The archived originals reproduce all nine optimized STLs byte for
byte.

Table~\ref{tab:bolts-meshes} records the individual smooth-shaft models.

\begin{table}[htbp]
\centering
\caption{Original and simplified copper bolt models; volume change is in percent.}
\label{tab:bolts-meshes}
\begin{tabular}{lrrrr}
\toprule
Bolt & Original faces & New faces & Shaft radius (mm) & Volume change (\%) \\
\midrule
1 & 1,690 & 192 & 8.5741 & 0.0000057 \\
2 & 1,690 & 192 & 8.5741 & 0.0000107 \\
3 & 1,690 & 192 & 8.5741 & 0.0000037 \\
4 & 2,224 & 192 & 8.3681 & 0.0000899 \\
5 & 1,690 & 192 & 8.5741 & 0.0000037 \\
6 & 1,690 & 192 & 8.5741 & 0.0000057 \\
7 & 1,690 & 192 & 8.5741 & 0.0000057 \\
8 & 2,224 & 192 & 8.3681 & 0.0000029 \\
9 & 1,690 & 192 & 8.5741 & 0.0000057 \\
\bottomrule
\end{tabular}
\end{table}

The optimized FTInnerShieldAndFlange STL is unchanged. Since shield and bolts share the same placement
transform, their collision test uses stored STL coordinates. VTK triangle collision finds zero surface
contacts, and enclosed-point checks in both directions find no bolt vertices inside the shield and no shield
vertices inside any bolt. Minimum bidirectional vertex-to-surface distances are 0.150002 mm for each bolt.
These sampled distances are not a certified global clearance bound; the independent surface and containment
checks address intersection. A ten-volume shield-plus-bolts Geant4 fixture completes without exceptions at
check\_overlaps=2 and again with a 10,000-point geometry test. The unchanged shield hash is retained with
the clearance results.

Tracking benchmarks isolate the nine copper bolts in air, excluding the shield from both timed fixtures.
Both cases use the new names, identical materials and placements, the same recorded GEMC3 executable,
FTFP\_BERT without optical processes, no magnetic field, one thread, default production cuts, and no GUI or
event output. Matched 5 GeV LUND primaries start at the target and sample solid angle uniformly from 1.3 to
3.0 degrees over full azimuth. Five repetitions with alternating fixture order and seeds 12345--12349 use
200 events per electron, proton, positive-pion, and neutron run. Time since beamOn measures transport;
separate zero-event runs measure startup. All 52 processes complete normally with all forty event counts
correct and no recorded Geant4 exceptions in either case.

Table~\ref{tab:bolts-timing} gives median transport wall seconds per 200 events.

\begin{table}[htbp]
\centering
\caption{Median transport wall seconds over five repetitions of the nine-bolt workload.}
\label{tab:bolts-timing}
\begin{tabular}{lrrr}
\toprule
Particle & Original & Optimized & Original/optimized \\
\midrule
e- & 0.551724 & 0.457388 & 1.206 \\
proton & 0.329016 & 0.329104 & 1.000 \\
pi+ & 0.330425 & 0.283490 & 1.166 \\
neutron & 0.133233 & 0.111502 & 1.195 \\
\bottomrule
\end{tabular}
\end{table}

Median startup simulation wall time is 0.659214 s originally and 0.623648 s after simplification. Peak
startup RSS falls from 129.34 to 115.31 MiB. The transport wall-time ratios are 1.206 for electrons, 1.000
for protons, 1.166 for positive pions, and 1.195 for neutrons. The proton result shows essentially unchanged
tracking time. These short workloads and five repetitions do not establish a universal speedup; material
redistribution changes stochastic shower histories even for matched primaries and seeds. Full CLAS12
performance and detector-response equivalence are not inferred.

Full default beamline construction completes without missing meshes and with no reported bolt-shield
overlap. Four remaining solid warnings and ten overlap reports concern other volumes, compared with four and
eleven immediately before this study. Sampled counts depend on tessellation and loading order, so this
difference is not claimed as a general overlap repair. All fifteen native geometry/material variations and
fourteen CAD variations pass the semantic comparator against GEMC2. Earlier apex, downstream,
adapter/exit-shield, adjuster, and inner-shield tracking measurements remain in their respective sections
with separate workloads and archives.

\FloatBarrier
Figure~\ref{fig:bolts-stl} compares a regular and a longer copper bolt.

\begin{figure}[htbp]
\centering
\includegraphics[width=\textwidth]{figures/bolts_stl_comparison.pdf}
\caption{Original threaded bolts (left) and smooth solid replacements (right).
Bolts 1 and 4 represent the regular and longer models; each pair shares bounds and camera settings.
Coordinates are recentered and permuted to put the bolt shaft on the horizontal axis.}
\label{fig:bolts-stl}
\end{figure}
